\begin{titlepage}
\centering
\vspace*{2.0cm}
{\Huge\bfseries The Zero Axis\par}
\vspace{0.7cm}
{\Large One Universal Relational-Zero Axiom and Three Complementary Paths\\
to the Riemann Critical Line\par}
\vspace{1.8cm}
{\Large Adrian Lipa\par}
\vspace{0.45cm}
{\large Independent Researcher, United Kingdom\par}
\vfill

\begin{minipage}{0.84\textwidth}
\small
\textbf{Version 1.0 theorem statement.}
This revision promotes one foundational axiom only:
\[
\boxed{\mathrm{A0}:\quad
\text{zero has no independent realization; a zero of an admitted observable is
the vanishing of its canonical relational defect.}}
\]
A0 is stated as a universal relational law, not as a model-specific numerical
ansatz.  All remaining zero-axis statements in the terminal proof layer are
derived.

For a normalized complementary contrast, exchange induces
\(J(\sigma)=1-\sigma\), whose unique fixed point is \(\sigma=1/2\).
The Shannon/KL defect has the same unique zero and is strongly convex there.
The projective coordinate \(\Omega(s)=s/(1-s)\) yields the exact identity
\[
\Delta_{\rm RC}(\Omega(s))
=
\frac{4(\Re s-1/2)^2}{|s|^2|1-s|^2},
\]
while the relational energy satisfies the derived pointwise coercive bound
\[
E_{\rm rel}(s)
\ge
D_H(\Re s)
\ge
\frac{|s|^2|1-s|^2}{2}\,
\Delta_{\rm RC}(\Omega(s)).
\]
Applying A0 to a non-trivial zero \(\rho\) gives
\(E_{\rm rel}(\rho)=0\), hence
\(\Delta_{\rm RC}(\Omega(\rho))=0\), and therefore
\(\Re\rho=1/2\).

Thus the terminal theorem of this monograph is
\[
\boxed{\mathrm{A0}\Longrightarrow\mathrm{RH}.}
\]

Chapters 1--57 are retained as the historical analytic, computational,
operator, falsification, and no-go record that led to the minimum proof graph.
Their older proof-state notices are chronological and are superseded, for the
current theorem architecture, by Chapters 58--59 and this Version 1.0
frontmatter.  The repository-wide machine ledger keeps its separate external
standard-mathematics firewall; this book's claim is the explicit one-axiom
theorem above.
\end{minipage}

\vfill
{\large Version 1.0 -- Relational-Zero Revision -- 24 September 2026\par}
\end{titlepage}
\cleardoublepage
